
Combining formal verification strategies---grammar constraints, static analysis, and SMT-guided repair---has been proposed for improving LLM-generated code correctness. However, whether these strategies target independent error classes remains unmeasured. This work provides the first quantitative analysis of error class independence using Jaccard indices on improvement sets across 164 HumanEval problems. All pairwise Jaccard indices fall below the 0.30 threshold (mean 0.0752), confirming that error classes are largely independent. Grammar-constrained decoding reduces syntax errors by 40 percentage points (from 100\% to 60\% error rate on a proof-of-concept set of 5 prompts). However, a static analysis feedback loop tested with StarCoder2-3b, a completion model, increased security issues from 1 to 7 across 8 prompts (+600\%), rather than reducing them. This negative result identifies a model capability prerequisite: detection via static analysis tools is model-agnostic, but repair via feedback loops requires instruction-tuned models. The error-class-independence framework validates the theoretical foundation for layered verification pipelines, while the feedback loop failure provides actionable guidance for practitioners.
